\section{Ресурсы}\label{ACME.Res}

\subsection{Перечень ресурсов}\label{ACME.Res.List}

Сервер ACME управляет следующими ресурсами: 
\begin{itemize}
\item
учетные записи (см.~\ref{ACME.Res.Account})
\item
заявки на выпуск сертификата (см.~\ref{ACME.Res.Order});
\item
объекты авторизации (см.~\ref{ACME.Res.Authz});   
\item
вызовы авторизации (см.~\ref{ACME.Res.Chall});
\item
выпущенные сертификаты (см.~\ref{ACME.Order.Cert});
\item
узел \code{directory} (см.~\ref{ACME.Res.Dir});
\item
узел \code{newNonce} (см.~\ref{ACME.Transport.Nonce});
\item
узел \code{newAccount} (см.~\ref{ACME.Account.Create});
\item
узел \code{newOrder} (см.~\ref{ACME.Order.Apply});
\item
узел \code{revokeCert} (см.~\ref{ACME.Revoke});
\item
узел \code{keyChange} (см.~\ref{ACME.Account.Rollover}).
\end{itemize}

Поддержка узлов \code{directory} и \code{newNonce} обязательна. 

Ресурсы описываются URL-адресами. Ресурсы могут ссылаться друг на друга, 
ссылки могут уточняться веб-отношениями~\cite{RFC8288}. 
%
Отношение \sstr{up} используется в ответе на URL-адрес вызова авторизации, 
указывая на адрес объекта авторизации, к которому вызов относится.
%
% skip: It is also used, with some media types, from certificate resources to
%   indicate a resource from which the client may fetch a chain of CA
%   certificates that could be used to validate the certificate in the original
%   resource.
%
Отношение \sstr{index} используется в ответах со всех узлов, кроме \code{directory}, 
указывая на адрес \code{directory}.

\begin{table}[htb]
\caption{Последовательность запросов клиента ACME}\label{Table.ACME.Seq}
\begin{tabular}{|l|p{5cm}|l|}
\hline
Действие & Запрос & Статус\\
\hline
\hline
Получить каталог & 
\code|GET| \code|directory| & 
200\\ 
\hline
%
Получить синхропосылку & 
\code|HEAD| \code|newNonce| & 
200\\
\hline
%
Создать учетную запись & 
\code|POST| \code|newAccount| & 
$201\rightarrow\text{учетная запись}$\\
\hline
%
Отправить заявку & 
\code|POST| \code|newOrder| & 
$201\rightarrow\text{заявка}$\\
\hline
%
Получить запросы авторизации & 
\code|POST-as-GET| адреса объектов авторизации в заявке &
200\\
\hline
%
Ответить на вызовы авторизации & 
\code|POST| адреса вызовов авторизации &
200\\
\hline
%
Получить статус обработки заявки & 
\code|POST-as-GET| заявка & 
200\\
\hline
%
Завершить формирование заявки & 
\code|POST| адрес завершения заявки & 
200\\
\hline
%
Получить статус обработки заявки & 
\code|POST-as-GET| заявка & 
200\\
\hline
%
Скачать сертификат & 
\code|POST-as-GET| адрес сертификата заявки & 
200\\
\hline
\end{tabular}
\end{table}

В таблице~\ref{Table.ACME.Seq} показана типичная последовательность запросов клиента, 
необходимых для создания новой учетной записи, подтверждения контроля над 
идентификатором, выпуска сертификата и его загрузки через некоторое время после 
выпуска. 
%
Стрелка вправо в таблице указывает на создаваемый ресурс. Ссылка на ресурс 
приводится в заголовке \code{Location} ответа. 

\subsection{Узел \code{directory}}\label{ACME.Res.Dir}

Для информирования клиентов об адресах ресурсов сервер поддерживает узел 
\code{directory}. В ответах с узла сервер высылает клиенту каталог ресурсов. 
Это объект JSON, поля которого описаны в таблице~\ref{Table.ACME.Dir}.

\begin{table}[b]
\caption{Каталог ресурсов ACME}\label{Table.ACME.Dir}
\begin{tabular}{|l|l|}
\hline
Поле & Значение (URL)\\
\hline
\hline
\sstr{newNonce} & Новая синхропосылка\\
%
\sstr{newAccount} & Новая учетная запись\\
%
\sstr{newOrder} & Новая заявка\\
%
\sstr{newAuthz} & Новая авторизация\\
%
\sstr{revokeCert} & Отзыв сертификата\\
%
\sstr{keyChange} & Смена ключа\\
\hline
\end{tabular}
\end{table}

URL-адрес каталога может быть произвольным при условии, что он отличается от 
адресов других узлов ACME и не конфликтует с адресами, не связанными с 
реализацией ACME.

% skip: For instance:
%  -  a host that functions as both an ACME and a Web server may want to
%     keep the root path "/" for an HTML "front page" and place the ACME
%     directory under the path "/acme".
%  -  a host that only functions as an ACME server could place the
%     directory under the path "/".

Если сервер ACME не реализует предварительную авторизацию
(см.~\ref{ACME.Res.Authz}), то он должен исключить из каталога поле
\sstr{newAuthz}.
 
Каталог может дополнительно содержать поле \sstr{meta} с метаданными о сервере
в виде объекта JSON. 
%
Поля объекта:

\begin{itemize}
\item
\sstr{termsOfService} (необязательное, строка)~--- URL-адрес с информацией об 
условиях обслуживания клиентов;

\item
\sstr{website} (необязательное, строка)~--- URL-адрес с дополнительной 
информацией о сервере;

\item
\sstr{caaIdentities} (необязательное, массив строк)~--- доменные имена, которые 
сервер считает относящимися к самому себе и которые можно использовать в 
записях авторизации УЦ (см.~\cite{RFC6844}).
%
Сервер должен задавать имена так, чтобы их можно было посимвольно 
переносить в поля \code{issue} и \code{issuewild} записей авторизации. 
%
Клиент может использовать имена для настройки своих записей авторизации;

\item
\sstr{externalAccountRequired} (необязательное, логическое значение)~--- 
наличие данного поля со значением \code{true} означает, что запросы на узел 
\code{newAccount} должны включать поле \sstr{externalAccountBinding}, которое 
связывает создаваемую учетную запись с внешней учетной записью. 
\end{itemize}

Клиент получает каталог ресурсов через GET-запрос на узел \code{directory}.

% skip: https://errata.rfc-editor.org/eid5771/
%   Before making a request to any URL from the directory, the client
%   MUST evaluate whether the directory object is still fresh according to
%   the Cache-Control header(s) received when that directory object was
%   accessed.  If no Cache-Control header(s) were received, the client MUST
%   act as if "Cache-Control: no-cache" was received.  If the directory
%   object is no longer fresh, the client MUST access the directory again
%   (by sending another GET request to the directory URL) and then use the
%   updated directory object.

\if0
HTTP/1.1 200 ОК
Content-Type: application/json

\begin{codeblock}
{
  "newNonce": "https://example.com/acme/new-nonce",
  "newAccount": "https://example.com/acme/new-account",
  "newOrder": "https://example.com/acme/new-order",
  "newAuthz": "https://example.com/acme/new-authz",
  "revokeCert": "https://example.com/acme/revoke-cert",
  "keyChange": "https://example.com/acme/key-change",
  "meta": {
    "termsOfService": "https://example.com/acme/terms/2017-5-30",
    "website": "https://www.example.com/",
    "caaIdentities": ["example.com"],
    "externalAccountRequired": false
  }
}
\end{codeblock}
\fi

\subsection{Учетные записи}\label{ACME.Res.Account}

С каждой учетной записью сервер связывает метаданные, которые описываются 
объектом JSON со следующими полями:

\begin{itemize}
\item
\sstr{status} (обязательное, строка)~--- статус учетной записи. 
Возможные значения: \sstr{valid} (действительна), \sstr{deactivated} 
(деактивирована) и \sstr{revoked} (отозвана).
%
Значение \sstr{deactivated} указывает на блокировку записи со стороны 
клиента, значение \sstr{revoked}~--- со стороны сервера 
(см.~\ref{ACME.Res.Status});

\item
\sstr{contact} (необязательное, массив строк)~--- набор URL-адресов для связи с 
клиентом по вопросам, которые касаются учетной записи, например, по вопросам 
отзыва или истечения срока действия сертификата. Поддерживаемые 
схемы URL-адресов определены в~\ref{ACME.Account};

\item
\sstr{termsOfServiceAgreeed} (необязательное, логическое значение)~--- 
согласие клиента с условиями обслуживания. Клиент подтверждает согласие, 
включая поле со значением \code{true} в запрос на узел \code{newAccount}.
Клиент не может обновлять поле;

\item
\sstr{externalAccountBinding} (необязательное, объект)~--- 
объект JWS, который привязывает внешнюю учетную запись клиента к учетной записи 
ACME (см.~\ref{ACME.Account.Binding}). Включая поле в запрос на узел 
\code{newAccount}, клиент подтверждает согласие на привязку. 
Клиент не может обновлять поле;

\item
\sstr{orders} (обязательное, строка)~--- URL-адрес, по которому размещен список 
связанных с учетной записью заявок на выпуск сертификатов. Список можно 
получить с помощью запроса POST-as-GET. 
\end{itemize}

\if0
\begin{codeblock}
{
  "status": "valid",
  "contact": [
    "mailto:cert-admin@example.org",
    "mailto:admin@example.org"
  ],
  "termsOfServiceAgreed": true,
  "orders": "https://example.com/acme/orders/rzGoeA"
}
\end{codeblock}
\fi

Список заявок в поле \sstr{orders} должен представлять собой массив строк,
каждая из которых~--- это URL-адрес отдельной заявки.
%
Серверу следует включать в список адреса заявок, которые ожидают  рассмотрения,
и не следует включать некорректные URL-адреса.
%
Сервер может возвращать неполный список, одновременно указывая в заголовке
\code{Link} ответа адрес, по которому можно получить продолжение списка, и
сопровождая этот адрес веб-отношением \sstr{next}.

\if0
\begin{codeblock}
HTTP/1.1 200 ОК
Content-Type: application/json
Link: <https://example.com/acme/directory>;rel="index"
Link: <https://example.com/acme/orders/rzGoeA?cursor=2>;rel="next"

{
  "orders": [
    "https://example.com/acme/order/TOlocE8rfgo",
    "https://example.com/acme/order/4E16bbL5iSw",
    "https://example.com/acme/order/neBHYLfw0mg"
  ]
}
\end{codeblock}
\fi

\subsection{Заявки на выпуск сертификата}\label{ACME.Res.Order}

Заявка на выпуск сертификата отражает запрос клиента на выпуск сертификата 
и используется для отслеживания хода обработки запроса. Заявка описывается 
объектом JSON со следующими полями:

% skip: Thus, the object contains information about the requested
%  certificate, the authorizations that the server requires the client
%  to complete, and any certificates that have resulted from this order.

\begin{itemize}
\item
\sstr{status} (обязательное, строка)~--- статус заявки.
%
Возможные значения: \sstr{pending} (ожидание рассмотрения), 
\sstr{ready} (готово), \sstr{processing} (обработка), 
\sstr{valid} (действительна) и \sstr{invalid} (недействительна) 
(см.~\ref{ACME.Res.Status});

\item
\sstr{expires} (необязательное, строка)~--- момент времени,
после наступления которого сервер будет считать заявку недействительной.
%
Время указывается в формате, заданном в~\cite{RFC3339}. 
%
Поле обязательно для объектов со статусами \sstr{pending} и \sstr{valid};

\item
\sstr{identifiers} (обязательное, массив объектов)~--- список объектов,
которые описывают идентификаторы заявки. 
%
Каждый объект включает два обязательных поля: \sstr{type} (строка)~---
тип идентификатора и \sstr{value} (строка)~--- собственно идентификатор.
%
В настоящем стандарте определяется только один тип идентификатора~--- 
\sstr{dns};

\item
\sstr{notBefore} (необязательное, строка)~--- отметка времени,
которая повторяет поле \code{notBefore} в запросе на выпуск 
сертификата. 
%
Время указывается в формате, заданном в~\cite{RFC3339};

\item
\sstr{notAfter} (необязательное, строка)~--- отметка времени,
которая повторяет поле \code{notAfter} в запросе на выпуск 
сертификата. 
%
Время указывается в формате, заданном в~\cite{RFC3339};

\item
\sstr{error} (необязательное, объект)~--- описание ошибки,
возникшей при обработке заявки, если таковая имеется. 
Описание представляет собой объект JSON с подробностями проблемы
(см.~\ref{ACME.Transport.Err});

\item
\sstr{authorizations} (обязательное, массив строк)~--- для заявки со статусом  
\sstr{pending}~--- авторизации, которые клиент должен выполнить перед 
выпуском запрошенного сертификата (см.~\ref{ACME.Authz}), включая уже 
выполненные клиентом неистекшие авторизации относительно идентификаторов 
заявки. 
%
Для заявки со статусом \sstr{valid} или \sstr{invalid}~---
выполненные авторизации. 
%
Строка массива \sstr{authorizations} представляет собой URL-адрес, 
по которому размещены адреса объекты авторизации. Объекты можно получить 
с помощью запроса POST-as-GET.

\begin{note*}
Перечень объектов авторизации определяется политикой сервера. Между 
идентификаторами и объектами может не быть соотношения 1:1;
\end{note*}

\item
\sstr{finalize} (обязательное, строка)~--- URL-адрес, на который необходимо 
отправить запрос на выпуск сертификата после завершения всех авторизаций.
%
При успешной обработке запроса будет сформирован адрес, по которому можно 
получить сертификат (см. следующее поле);

\item
\sstr{certificate} (необязательное, строка)~--- URL-адрес, по которому 
размещен сертификат, выпущенный в результате обработки заявки. 
\end{itemize}

\if0
\begin{codeblock}
{
  "status": "valid",
  "expires": "2026-01-20T14:09:07.99Z",

  "identifiers": [
    { "type": "dns", "value": "www.example.org" },
    { "type": "dns", "value": "example.org" }
  ],

  "notBefore": "2026-01-01T00:00:00Z",
  "notAfter": "2026-01-08T00:00:00Z",

  "authorizations": [
    "https://example.com/acme/authz/PAniVnsZcis",
    "https://example.com/acme/authz/r4HqLzrSrpI"
  ],

  "finalize": "https://example.com/acme/order/TOlocE8rfgo/finalize",

  "certificate": "https://example.com/acme/cert/mAt3xBGaobw"
}
\end{codeblock}
\fi
 
В заявке, которая отправляется на узел \code{newOrder}, идентификатор 
типа \sstr{dns} может представлять собой подстановочное доменное имя. 
%
Такое имя состоит из одного символа звездочки, за которым следует один символ
точки (\sstr{*.}), а затем идет обычное доменное имя,  заданное в соответствии с
правилами СТБ~34.101.19 (см. расширение \code{SubjectAltName}).
%
В объекте авторизации, который сервер формирует для подстановочного доменного имени, 
идентификатор не должен включать префикс \sstr{*.}. Объект должен включать 
необязательное поле \sstr{wildcard} со значением \code{true}. 

После формирования массивов \sstr{identifiers} и \sstr{authorizations} 
они не должны изменяться. Если клиент обнаруживает изменения, то он 
должен считать заявку недействительной.  

В массиве \sstr{authorizations} следует указывать все объекты авторизации, 
которые УЦ учитывает при принятии решения о выпуске сертификата. В том числе те
объекты, которые использовались при обработке предыдущих заявок или на этапе
предварительной авторизации.
%
% skip: For example, if a CA allows multiple
%  orders to be fulfilled based on a single authorization transaction,
%  then it SHOULD reflect that authorization in all of the orders.

Соответственно наличие URL-адреса авторизации в массиве \sstr{authorizations} 
не означает автоматически, что клиент обязан выполнить какие-либо действия:
клиент мог выполнить авторизацию в рамках обработки предыдущей заявки, 
предварительно (см.~\ref{ACME.Authz}) или сервер мог авторизовать клиента 
на основе внешней учетной записи.   
%
Клиенту следует проверять поле \sstr{status} заявки, чтобы определить, нужно ли 
ему предпринимать действия по авторизации.

\subsection{Объекты авторизации}\label{ACME.Res.Authz}
 
Авторизация в ACME подтверждает право владения сертификатами с определенным 
идентификатором со стороны клиента с определенной учетной записью.
%
% info: авторизация --- назначение прав доступа (Bias)
%
Авторизация описывается объектом JSON, который помимо идентификатора включает
метаданные (например, статус авторизации) и перечень вызовов для подтверждения
контроля над идентификатором со стороны владельца учетной записи.

Объект авторизации состоит из следующих полей:
%
\begin{itemize}
\item
\sstr{identifier} (обязательное, объект)~--- объект, который описывает целевой 
идентификатор.
%
Объект включает два обязательных поля: \sstr{type} (строка)~---
тип идентификатора и \sstr{value} (строка)~--- собственно идентификатор;

\item
\sstr{status} (обязательное, строка)~--- статус авторизации.
%
Возможные значения: \sstr{pending} (ожидание рассмотрения), 
\sstr{valid} (действительна), \sstr{invalid} (недействительна),
\sstr{deactivated} (деактивирована), \sstr{expired} (истек срок действия)
и \sstr{revoked} (отозвана) (см.~\ref{ACME.Res.Status});

\item
\sstr{expires} (необязательное, строка)~--- момент времени,
после наступления которого сервер будет считать авторизацию недействительной. 
%
Время указывается в формате, заданном в~\cite{RFC3339}. 
%
Поле обязательно для объектов со статусом \sstr{valid}; 

\item
\sstr{challenges} (обязательное, массив объектов)~--- для объекта авторизации 
со статусом \sstr{pending}~--- список вызовов авторизации, которые докажут 
контроль над идентификатором.
%
Клиенту следует отработать один из вызовов, а серверу следует считать контроль 
доказанным, если один из вызовов успешно отработан.
%
Для объекта авторизации со статусом \sstr{valid} массив \sstr{challenges} 
содержит успешно отработанный вызов, а для объекта со статусом 
\sstr{invalid}~--- вызов, который был отработан с ошибкой;

\item
\sstr{wildcard} (необязательное, логическое значение)~--- поле должно
присутствовать и принимать значение \code{true}, если объект авторизации
создан в ответ на заявку на узел \code{newOrder}, в которой в качестве 
идентификатора указано подстановочное доменное имя (см.~\ref{ACME.Res.Order}).
%
В других случаях поле должно или отсутствовать, или принимать значение 
\code{false}.
%
% skip: For pre-authorizations, it MUST be absent or false. Wildcard domain 
% names are described in Section 7.1.3.
% use: https://errata.rfc-editor.org/eid6364/.
\end{itemize}

В настоящем стандарте поддерживаются только идентификаторы типа \sstr{dns},
т.~е. полные доменные имена. Доменное имя должно быть представлено в той же
форме, в которой оно будет указано в сертификате, т.~е. в соответствии с 
правилами~СТБ~34.101.19 (подраздел 9.3).
%
Сервер должен проверять корректность кодового представления 
интернационализированных доменных имен (с префиксом \sstr{xn--} в соответствии  
с~\cite{RFC5890}).

Подстановочное доменное имя (со звездочкой) не должно включаться в объект 
авторизации. Даже если заявка на узел \code{newOrder} содержит подстановочное
доменное имя, сервер должен включить в объект авторизации полное доменное имя,
сопроводив его значением \code{true} в поле \sstr{widlcard}.

\if0
\begin{codeblock}
{
  "status": "valid",
  "expires": "2025-03-01T14:09:07.99Z",

  "identifier": {
    "type": "dns",
    "value": "www.example.org"
  },

  "challenges": [
    {
      "url": "https://example.com/acme/chall/prV_B7yEyA4",
      "type": "http-01",
      "status": "valid",
      "token": "DGyRejmCefe7v4NfDGDKfA",
      "validated": "2024-12-01T12:05:58.16Z"
    }
  ],

  "wildcard": false
}
\end{codeblock}
\fi

\subsection{Вызовы авторизации}\label{ACME.Res.Chall}

Вызов авторизации~--- это объект JSON, через который сервер предлагает 
клиенту подтвердить контроль над определенным идентификатором.
%
Структура вызова определяется методом проверки контроля и поэтому 
вызов не имеет строго определенного формата.
%
Общая структура вызова, устанавливаемая настоящим стандартом, методы проверки 
и соответствующие форматы определяются в разделе~\ref{ACME.Chall}.

\subsection{Состояния объектов}\label{ACME.Res.Status}

Объект ACME проходит несколько стадий обработки. Состояние объекта~--- стадия 
его обработки~--- указывается в поле \sstr{status}.

Вызов авторизации создается в состоянии \sstr{pending}. Он переходит в 
состояние \sstr{processing}, когда клиент отрабатывает вызов 
(см.~\ref{ACME.Authz}) и сервер начинает проверять, завершена ли обработка.
%
Если проверка прошла успешно, вызов переводится в состояние~\sstr{valid}, в 
случае ошибки~--- в состояние \sstr{invalid}.
%
Состояние \sstr{pending} не изменяется, когда сервер повторно проверяет 
отработку вызова (см.~\ref{ACME.Chall}) и когда клиент повторно отрабатывает 
вызов.

Объект авторизации создается в состоянии \sstr{pending}. Если один из вызовов, 
указанных в объекте, переходит в состояние \sstr{valid}, объект также переходит 
в состояние \sstr{valid}. Если клиент выполняет отработку вызова с ошибкой
или если происходят другие ошибки, а объект все еще в состоянии \sstr{pending},
состояние меняется на \sstr{invalid}. Состояние \sstr{valid}
может измениться на \sstr{expires} по истечении срока действия объекта,
\sstr{deactivated} при деактивации объекта клиентом и на 
\sstr{revoked} при отзыве объекта сервером.

Заявки создаются в состоянии \sstr{pending}. Как только все объекты 
авторизации, указанные в заявке, переходят в состояние \sstr{valid}, заявка 
переходит в состояние \sstr{ready}. Заявка переходит в состояние 
\sstr{processing} после того, как клиент отправляет запрос на URL-адрес 
финализации заявки, и УЦ начинает выпуск сертификата. 
%
После выпуска сертификата заявка переходит в состояние \sstr{valid}. 
%
Если на любом из этапов возникает ошибка, заявка переходит в состояние 
\sstr{invalid}. В это состояние заявка переходит также тогда, когда 
истек срок ее действия и состояние одного из объектов авторизации изменилось
с \sstr{valid} на \sstr{expires}, \sstr{revoked} или \sstr{deactivated}).

Учетные записи создаются в состоянии \sstr{valid}, поскольку после 
успешного запроса на узел \code{newAccount} никаких дополнительных действий для 
создания учетной записи не требуется. Если учетная запись деактивируется 
клиентом или отзывается сервером, она переводится в соответствующее состояние.

\begin{note*}
Сервер может обрабатывать объекты так, что они не попадают в некоторые из 
перечисленных состояний. Например, сервер может обрабатывать заявку <<на лету>>
и тогда она не попадет в состояние \sstr{processing}. Или сервер может удалять
объект авторизации как только истек срок его действия, и тогда объект не 
попадает в состояние \sstr{expired}.
\end{note*}

